\section{Stability of the explicit step and step bounds}\label{sup:steps}

% The paper cites these subsections by number (S2.1 to S2.4): keep their order.

\subsection{Coupling with the engines}

The contact model knows nothing of the engine: bodies are indices, and the friction state is
passed in and out so that each engine can keep it where its stepping scheme requires. The
backends are thin: each one builds the engine's model of a scenario, and at every step reads the
state, evaluates the contact model, applies the wrenches as external forces and advances. The
friction state advances once per step, from the same evaluation. The contact models of the
engines themselves are switched off for the methods of the library.
\begin{itemize}
\item \emph{Drake} (version 1.58). The bodies are free bodies of a discrete
  \code{MultibodyPlant} without contact geometry. A system connected to the plant outputs the
  wrenches as externally applied spatial forces and holds the friction state as discrete state
  with the period of the plant, and serves to the plant the evaluation made at the start of the
  step, so that the forces act on the state they were computed from.
\item \emph{Project Chrono} (version 10.0). The wrenches are accumulated on the bodies at their
  centres of mass before each step of the linearized implicit Euler time stepper, for which an
  external load is a constant.
\item \emph{MuJoCo} (version 4.0). Bodies have free joints; the wrenches are written as applied
  Cartesian forces and the step is taken with the \code{implicitfast} integrator, whose implicit
  terms do not include applied forces.
\end{itemize}
Each backend also passes its step and the inverse inertias of the bodies to the contact model,
which is what the bounds of Section~\ref{sup:bounds} need; the engines store the state of a
pair as 26 numbers. For a free body under an applied force the three steppers are the same
first-order scheme, the semi-implicit Euler step: the velocity is advanced with the force of
the start of the step and the position with the new velocity. They differ in how they integrate
rotation and gyroscopic terms.

\subsection{Stability condition}

The step is $\vv^+ = \vv + \dt\,M^{-1}\vect{f}(\x,\vv)$, $\x^+ = \x + \dt\,\vv^+$. For a spring
$k$ and a damper $c$ acting on an inverse mass $w$, with $\alpha = k\,w\,\dt^2$ and
$\beta = c\,w\,\dt$, one step maps $(x, \dt\,v)$ with the matrix
$\bigl(\begin{smallmatrix} 1-\alpha & 1-\beta\\ -\alpha & 1-\beta\end{smallmatrix}\bigr)$, whose characteristic
polynomial is $\lambda^2 - (2-\alpha-\beta)\lambda + (1-\beta)$. By the Jury conditions the scheme is stable
if and only if $0 < \alpha$, $\beta < 2$ and $\alpha + 2\beta < 4$ (at $\alpha = 0$ it is only marginally stable), that is,
\begin{equation}
  k\,w\,\dt^2 + 2\,c\,w\,\dt < 4 \quad\text{and}\quad c\,w\,\dt < 2 .
  \label{eq:sup-stability}
\end{equation}
When several contacts act on a body their contributions add. Four mechanisms of the laws enter
this condition, and three of them are proportional to the load: the stick spring and damper;
the normal damper, $c = h\,K V$; and the sliding friction across the stick--slip transition,
which is a damper of up to $0.86\,\mu\Fn/v_s$. For a body of mass $m$ under a contact load of
$N$ times its weight, with $h = 5$\,s/m, $\mu = 0.5$ and $\dt = 1$\,ms, the stick has
$\alpha + 2\beta = 0.54\,N$, so the default bristle is stable up to a load of about seven times
the weight; the normal damper has $\beta = 0.049\,N$ and the transition $\beta = 0.17\,N$. A
body under its own weight is far from all three limits. A block at the bottom of a tall stack,
or a light body under a heavy one, is not.

The fourth mechanism is the normal stiffness, which is the model's and is not touched, since
changing it would change the statics. For a flat contact of area $A_c$ between a body of mass
$m$ and the ground it needs $K A_c\,\dt^2/m < 4$; for a body held between two such contacts the
two stiffnesses add, $\sum w = 2/m$, and with the dampers at their bounds the condition
$K A_c\,\dt^2\sum w < 2$ gives $\dt < \sqrt{m/(K A_c)}$. The stiffnesses of the battery were
chosen so that the base step of 1\,ms satisfies it, with the two exceptions that are there to
show the limit.

\subsection{Bounds on the terms proportional to the load}\label{sup:bounds}

Let $\Lambda = \sum|\Fn|\,w$ be the \emph{load acceleration} of a body, the sum over its
contacts of the normal force times the inverse mass that the contact sees. With $\Lambda$ the
sum of the load accelerations of the two bodies of a pair, each pair applies
\begin{equation}
  k \le \frac{|\Fn|}{\Lambda\,\dt^2}, \qquad c \le \frac{|\Fn|}{\Lambda\,\dt}, \qquad
  h \le \frac{1}{\Lambda\,\dt}, \qquad v_s \ge \mu\,\Lambda\,\dt ,
  \label{eq:bounds}
\end{equation}
and the torsional bristle the same with the inverse inertia about the normal, so that each
mechanism satisfies $\sum k\,w\,\dt^2\le1$ and $\sum c\,w\,\dt\le1$ over the contacts of a body,
which keeps Eq.~\eqref{eq:sup-stability} with margin when two mechanisms act on one body. In
words: a stick spring cannot bring the pair back in less than a step, a damper cannot take more
than the relative velocity of a step, and the stick--slip transition cannot be narrower than
the change of velocity that friction causes in a step. The load acceleration counts the lever
arm of each contact and, for the normal damper, the rocking that the damping of the foundation
resists; it is at least the pair's own. $\Lambda$ is taken from the previous evaluation and
kept in the state of the pair. The bounds are per unit of load, so they do not depend on how a
contact is cut into pieces, and there is no sub-stepping. The baselines of the library are
subject to the same bounds.

\subsection{What the bounds change}

Below the threshold of Eq.~\eqref{eq:bounds} nothing changes: a body under its own weight at
1\,ms receives the same wrench to the last bit with and without the bounds, which a test of the
library asserts. Above it, the parameters that the user gave are no longer the ones applied: the
pre-sliding displacement grows from $d$ to $\mu\Lambda\,\dt^2$, the stick--slip transition widens
to $\mu\Lambda\,\dt$, and the damping factor falls below $h$, all as functions of the step and of
the load; the dampers that a light body carries are bounded by its mass over the step, so a
heavy body resting on a light one is damped less than the model prescribes. What is benchmarked
there is the law regularized by the step, and its results converge to the model's only as the
step is reduced. The battery did not record in which runs the bounds were active; from the
loads, they are in \scn{tower}, \scn{massratio} and \scn{stack} at the base step and are not in
the scenarios with one body under its own weight.
